\section{Public Power}
\label{sec:publicpower}

\subsection{Motivation}

At first sight, Public Power appears to be the sector for which no real
downscaling is needed. Public electricity and heat production is a classical
\emph{point-source} activity: emissions originate from plants whose locations
are physically defined. One would therefore expect the inventory coordinates to
be sufficient, so that no additional spatial redistribution would be needed.

Inspection of the CAMS-REG-ANT GNFR~A inventory and of the CAMS methodology
showed that this assumption is only partly true. Two observations justify a
dedicated proxy:
\begin{itemize}
  \item GNFR~A contains both point sources and area sources, so the sector is
        not represented exclusively by explicit stacks.
  \item Even where CAMS labels a source as a point, the reported coordinates
        are not always precise enough for high-resolution allocation.
\end{itemize}

The area component is not an accident of the file structure. In the CAMS
compilation logic, national totals for public electricity and heat production
are compared against the emissions that can be assigned to known plants. The
remaining difference is retained as residual \emph{area-source} mass and is
distributed spatially using population as a first-order allocator
\cite{CAMS_glob_ant_61}. Public Power is therefore a mixed spatial system: it
is dominated by point sources, but it also contains a residual area component
produced by the inventory construction itself.

The upgraded workflow therefore creates two proxy layers:
\begin{itemize}
  \item a point-source localisation layer for the explicit plant component;
  \item an area-source downscaling layer for the residual surface emissions.
\end{itemize}

\subsection{Point-source proxy: localisation of combustion units}

\paragraph{Rationale.}
For true power-plant emissions, no classical downscaling is required in the
sense of spreading emissions over land-cover classes. What is required is a
better localisation layer than the raw CAMS point coordinates. The upgraded
method therefore treats point-source Public Power as a \emph{localisation}
problem: identify the pixels that plausibly host thermal generation units and
use that support to anchor the point component.

The point layer is built from a European open power-plant register described
in Appendix~\ref{app:powerplant_register}. Only combustion-based public-power
units are retained, while hydro and other non-combustion generation types are
excluded because the objective is to represent stack-related emissions.

\paragraph{Proxy definition.}
Let $g$ denote a pixel on the national CORINE reference grid and let
$U_{\mathrm{comb}}$ be the set of retained combustion units. The point proxy is
a binary support raster:
\begin{equation}
  \mu_{g}^{\mathrm{pp,point}} =
  \begin{cases}
    1, & \exists\,u \in U_{\mathrm{comb}} \text{ with } x_u \in g,\\
    0, & \text{otherwise},
  \end{cases}
\end{equation}
where $x_u$ is the location of unit $u$. In other words, a reference-grid
pixel receives value one if at least one retained combustion unit falls inside
it. The binary choice is deliberate: it expresses plant presence without
letting a few large multi-unit sites dominate purely through registry
granularity.

\paragraph{Interpretation.}
This layer should not be read as a full reconstruction of plant-level
emissions. It is a spatial support surface answering a narrower question:
where are public-power point emissions most plausibly located at fine
resolution?

\subsection{Area-source proxy: residual CAMS mass distributed over suitable land}

\paragraph{Rationale.}
The existence of CAMS GNFR~A area emissions means that a second proxy is
necessary. Here a genuine downscaling problem does arise, because the residual
public-power mass must be redistributed from coarse CAMS area cells to fine
pixels. The chosen proxy follows the logic already embedded in the CAMS
methodology: if the residual originates from the gap between national totals
and explicitly reported plants, and CAMS distributes that gap using
population, then a population-based area proxy is methodologically coherent
rather than arbitrary.

The area layer combines three ingredients: the CAMS area-source geometry, a
CORINE-based suitability mask, and a gridded population field that controls
the within-cell intensity gradient.

\paragraph{Proxy definition.}
Let $k$ denote a CAMS GNFR~A area cell and let $g \in k$ index fine pixels
inside that cell on the reference grid. Define the CORINE eligibility flag
\begin{equation}
  I_g^{\mathrm{CLC}} =
  \begin{cases}
    1, & \mathrm{CLC}_g \in \mathcal{C}_{\mathrm{pp}},\\
    0, & \text{otherwise},
  \end{cases}
\end{equation}
and let $P_g$ be the LandScan population value resampled to pixel $g$. The raw
within-cell score is
\begin{equation}
  \mu_{g,k}^{\mathrm{pp,area}} =
  I_g^{\mathrm{CLC}} \cdot \max(P_g, P_0)^{\beta},
\end{equation}
where $P_g$ is the population value in pixel $g$. The
normalised within-cell weight is then
\begin{equation}
  w_{g,k}^{\mathrm{pp,area}} =
  \frac{\mu_{g,k}^{\mathrm{pp,area}}}
       {\displaystyle\sum_{g' \in k} \mu_{g',k}^{\mathrm{pp,area}}},
\end{equation}
so that the weights inside each CAMS area cell sum to unity.

This separates two roles clearly: land cover defines \emph{where} residual
public-power emissions may be placed, while population defines \emph{how much}
of the CAMS cell total each eligible pixel receives.

\paragraph{Fallback logic.}
Some CAMS area cells contain no pixels belonging to the selected CORINE
classes. In that case, the implementation falls back first to
population-only weighting over all valid pixels in the cell and, if necessary,
to a uniform distribution over the valid cell footprint. This avoids losing
area-source mass purely because the high-resolution land-cover layer does not
intersect the coarse CAMS rectangle in the expected way.

\subsection{How the point and area layers are used together}

GNFR~A is represented by two complementary spatial products because the
inventory itself contains two geometries. For area sources, pollutant $p$ from
CAMS source $s$ is redistributed within its parent CAMS cell using
\begin{equation}
  e_{g,s,p}^{\mathrm{area}} = E_{s,p} \, w_{g,s}^{\mathrm{pp,area}}.
\end{equation}
For point sources, the role of the proxy is different: emissions are not
spread across land-cover classes but constrained to pixels for which
\(\mu_{g}^{\mathrm{pp,point}} > 0\), that is, pixels hosting at least one
retained combustion unit.

Public Power is therefore not treated as a single undifferentiated proxy.
Instead, the method does two things at once: it \emph{relocates} the point
component more credibly, and it \emph{downscales} the residual area component
in a way consistent with the CAMS methodology.

\subsection{Limitations and forward path}

The main limitations are straightforward. The point layer is a binary plant
presence raster rather than a one-to-one reconciliation of individual plant
records. The plant register is also incomplete, so omission and location error
remain possible. Finally, the area layer assumes that residual public-power
emissions follow population once restricted to suitable land. That assumption
is defensible because it mirrors the inventory logic, but it does not resolve
the internal diversity of diffuse public-power activities.
